\section{Stacks of pair registers and the elimination of pair effects}
\label{sec:party_effect_stacks}

This section proves the statements about parties in the elimination of
normalized pair effects \cite[Lemma~5.1]{OpenAI2026PolynomialPEPS}. Every effect
occurrence of a monomial receives its own stack of $m$ pair registers, owned by
the two endpoint parties of the effect; a circuit of pair sources and pair
effects is rewritten, by one recursion on the circuit, as an effect chain whose
replaced terms are words of local contractions and normalized pair sources. For a
weighted gate this gives the elimination theorem on parties, all clauses of the
manuscript's lemma except the combination of sources on one pair of parties.

Throughout, $P$ is a set of parties, a system $X$ is a family of finite basis
sets $X_s$, $s\in P$, with configurations $X_P=\prod_sX_s$, and the operator of a
circuit $c$ from $X$ to $Y$ is the matrix $\mathrm{op}(c)\in\C^{Y_P\times X_P}$
of Definition~\ref{def:party_circuit}. The circuit $c$ followed by $c'$ is
written $c'\circ c$, so $\mathrm{op}(c'\circ c)=\mathrm{op}(c')\,\mathrm{op}(c)$.
For distinct parties $p\ne q$ and finite sets $A,B$, the system $T^{pq}_{A,B}$ has
the basis $A\times T_p$ at $p$, the basis $B\times T_q$ at $q$ and $T_s$ at every
other party (Definition~\ref{def:party_add_pair}); its configurations are
identified with $(A\times B)\times T_P$. For a finite set $\alpha$ and $n\in\N$
write $\alpha^n$ for the maps $\Fin{n}\to\alpha$.

\subsection{Stacks of pair registers}
\label{sec:party_stacks}

\begin{definition}[Circuits beside fresh factors at two parties]
    \label{def:party_lift_pair}
    \lean{QuantumCircuit.Party.PartyCircuit.liftPair,
        QuantumCircuit.Party.PartyCircuit.isAllowed_liftPair,
        QuantumCircuit.Party.PartyCircuit.gateCells_liftPair,
        QuantumCircuit.Party.PartyCircuit.gateCount_liftPair,
        QuantumCircuit.Party.Word.framePairs,
        QuantumCircuit.Party.Word.isAllowed_framePairs,
        QuantumCircuit.Party.Word.sources_framePairs,
        QuantumCircuit.Party.Word.gateCount_framePairs}
    \leanok
    \uses{def:party_circuit, def:party_add_pair, def:party_gate_frame,
        def:party_relabel, def:party_allowed, def:party_gate_counts, def:party_word}
    Let $p\ne q$ and let $A,B$ be finite sets. For a circuit $c$ from $S$ to $T$, the
    circuit $c$ beside the fresh factors $A$ at $p$ and $B$ at $q$ is the circuit
    from $S^{pq}_{A,B}$ to $T^{pq}_{A,B}$ that reads every basis of $S^{pq}_{A,B}$
    as the product of the factor $A$, $B$ or a one-point set with the basis of
    $S$, runs $c$ framed by these factors, and reads the bases back. It has the gate
    cells and the gate count of $c$. For an interpretation compatible with framing,
    with one outcome per gate, and a condition on gates that passes from a gate to
    its framed gates, it is allowed when $c$ is.

    For $n\in\N$ and a word $w$ from $S$ to $T$, the word beside a stack of $n$ pair
    registers on $p,q$ with pair basis $\alpha\times\beta$ is $w$ beside the fresh
    factors $\alpha^n$ at $p$ and $\beta^n$ at $q$. It is allowed when $w$ is, and
    it has the pair sources and the gate count of $w$.
\end{definition}

\begin{theorem}[Operators beside fresh factors]
    \label{thm:party_op_lift_pair}
    \lean{QuantumCircuit.Party.PartyCircuit.op_liftPair,
        QuantumCircuit.Party.Word.op_framePairs,
        QuantumCircuit.Party.Word.op_liftPair_mul_op_sourceAt}
    \leanok
    \uses{def:party_lift_pair, def:party_stack_system, def:party_pair_gates}
    For an interpretation compatible with framing, with one outcome per gate, the
    operator of a circuit $c$ beside the fresh factors $A$ at $p$ and $B$ at $q$ is
    $1_{A\times B}\otimes\mathrm{op}(c)$, read along the identifications of the
    configurations of $S^{pq}_{A,B}$ and $T^{pq}_{A,B}$ with $(A\times B)\times S_P$
    and $(A\times B)\times T_P$. In particular a word beside a stack of $n$ pair
    registers has the operator $1_{(\alpha\times\beta)^n}\otimes\mathrm{op}(w)$, the
    stack configurations read as in Definition~\ref{def:party_stack_system}. If
    $\sigma_\eta^S$ denotes the pair source preparing
    $\eta\in\C^{\alpha\times\beta}$ on fresh factors in front of the bases of $p$ and
    $q$ of $S$, then for every word $c$ from $S$ to $T$
    \begin{align}
        \mathrm{op}(\tilde c)\,\mathrm{op}(\sigma_\eta^S)
         & =\mathrm{op}(\sigma_\eta^T)\,\mathrm{op}(c),
        \notag
    \end{align}
    where $\tilde c$ is $c$ beside the fresh factors $\alpha$ at $p$ and $\beta$ at
    $q$ \cite[Section~5]{OpenAI2026PolynomialPEPS}.
\end{theorem}

\begin{proof}\leanok
    \uses{thm:party_op_frame, def:party_relabel, thm:party_word_op}
    The operator of a composition is the product of the operators; the two
    relabellings are the permutation matrices of the identifications of bases, and
    the framed circuit has the operator $1\otimes\mathrm{op}(c)$ by
    Theorem~\ref{thm:party_op_frame}. Comparing entries gives the first identity,
    and the second is its case $A=\alpha^n$, $B=\beta^n$. The operator of the pair
    source is the preparation $x\mapsto\eta\otimes x$; both sides of the last
    identity have the entry $\eta_{(a,b)}\,\mathrm{op}(c)_{yx}$ at the output
    configuration $((a,b),y)$ and the input $x$.
\end{proof}

\begin{definition}[Stacks of pair registers]
    \label{def:party_stack_system}
    \lean{QuantumCircuit.Party.stackSys, QuantumCircuit.Party.stackSys_of_ne,
        QuantumCircuit.Party.stackConfigEquiv, QuantumCircuit.Party.stackZeroEquiv,
        QuantumCircuit.Party.pairSourceOf, QuantumCircuit.Party.stackSource}
    \leanok
    \uses{def:party_add_pair, def:party_pair_gates, def:sites_config}
    Let $p\ne q$, let $\alpha,\beta$ be finite sets and $n\in\N$. The system $T$ with
    a stack of $n$ pair registers on $p,q$ is $T^{pq}_{\alpha^n,\beta^n}$; at every
    party other than $p$ and $q$ it has the basis of $T$. Its configurations are
    identified with $(\alpha\times\beta)^n\times T_P$, the stack configuration
    $j\mapsto(a_j,b_j)$ collecting the factor $a\in\alpha^n$ at $p$ and
    $b\in\beta^n$ at $q$. For $n=0$ every basis of the stacked system is identified
    with the corresponding basis of $T$. For $\eta\in\C^{\alpha\times\beta}$, the
    pair source with endpoints $p,q$, fresh factors $\alpha,\beta$ and vector $\eta$
    is the occurrence $(p,q,\eta)$; its gate from $T^{pq}_{\alpha^n,\beta^n}$ to
    $T^{pq}_{\alpha^{n+1},\beta^{n+1}}$ on the cell $\{p,q\}$ prepares $\eta$ at
    position $0$ of the stack, reading $a\in\alpha^{n+1}$ as
    $(a_0,(a_1,\dots,a_n))$ and likewise at $q$ \cite[Lemma~5.1]{OpenAI2026PolynomialPEPS}.
\end{definition}

\begin{definition}[Preparation of a stack]
    \label{def:party_prep_word}
    \lean{QuantumCircuit.Party.Word.prepWord,
        QuantumCircuit.Party.Word.isAllowed_prepWord,
        QuantumCircuit.Party.Word.sources_prepWord,
        QuantumCircuit.Party.Word.sourceCount_prepWord,
        QuantumCircuit.Party.Word.usesOnly_prepWord}
    \leanok
    \uses{def:party_stack_system, def:party_word, def:party_relabel}
    The preparation of $n$ copies of $\eta\in\C^{\alpha\times\beta}$ on $p,q$ is the
    word from $T$ to $T^{pq}_{\alpha^n,\beta^n}$ defined by induction on $n$: for
    $n=0$ it is the relabelling of the bases, and for $n+1$ it is the preparation of
    $n$ copies followed by the pair source preparing $\eta$ at position $0$. Its pair
    sources are $n$ copies of $(p,q,\eta)$, so it has $n$ sources; it uses only the
    parties $p$ and $q$, and it is allowed when $\lVert\eta\rVert=1$.
\end{definition}

\begin{theorem}[The prepared stack]
    \label{thm:party_op_prep_word}
    \lean{QuantumCircuit.Party.Word.op_prepWord}
    \leanok
    \uses{def:party_prep_word, def:party_effect_matrices}
    Read along the identification of the configurations of
    $T^{pq}_{\alpha^n,\beta^n}$ with $(\alpha\times\beta)^n\times T_P$, the
    preparation of $n$ copies of $\eta$ has the operator
    $x\mapsto\eta^{\otimes n}\otimes x$
    \cite[Lemma~5.1]{OpenAI2026PolynomialPEPS}.
\end{theorem}

\begin{proof}\leanok
    \uses{def:party_relabel, thm:party_word_op, def:party_effect_matrices}
    By induction on $n$. For $n=0$ the relabelling has the permutation matrix of
    the identification of bases. For $n+1$, the operator of the pair source is the
    preparation of $\eta$ in front of its input, and the entries of the product
    match since
    $\eta^{\otimes(n+1)}(f)=\eta(f_0)\,\eta^{\otimes n}(f_1,\dots,f_n)$.
\end{proof}

\begin{definition}[Insertion words]
    \label{def:party_insertion_word}
    \lean{QuantumCircuit.Party.Word.insEquiv, QuantumCircuit.Party.Word.insWord,
        QuantumCircuit.Party.Word.isAllowed_insWord,
        QuantumCircuit.Party.Word.sources_insWord,
        QuantumCircuit.Party.Word.usesOnly_insWord}
    \leanok
    \uses{def:party_prep_word, def:party_pair_gates, def:party_relabel}
    Let $g$ be a pair effect from $Z$ to $Y$ on the endpoints $p\ne q$, with pair
    basis $\alpha\times\beta$ and vector $\eta$, so that $Z_p=\alpha\times Y_p$ and
    $Z_q=\beta\times Y_q$ up to the identifications of $g$. For $m\ge1$ and
    $k\in\Fin{m}$, the insertion word of $g$ at position $k$ is the word from $Z$ to
    $Y^{pq}_{\alpha^m,\beta^m}$ that prepares $m-1$ copies of $\eta$ on $p,q$ and
    then relabels the bases: at $p$ the basis
    $\alpha^{m-1}\times(\alpha\times Y_p)$ becomes $\alpha^m\times Y_p$, the factor
    $\alpha$ removed by $g$ going to position $k$ of the stack and the prepared
    copies to the other positions in order, and likewise at $q$. It is allowed when
    $\lVert\eta\rVert=1$, its pair sources are $m-1$ copies of $(p,q,\eta)$, and it
    uses only $p$ and $q$ \cite[Lemma~5.1]{OpenAI2026PolynomialPEPS}.
\end{definition}

\begin{theorem}[The insertion word is the insertion map]
    \label{thm:party_op_insertion_word}
    \lean{QuantumCircuit.Party.Word.op_insWord}
    \leanok
    \uses{def:party_insertion_word, def:party_insertion_matrices,
        def:peps_cyclic_insertion}
    Read along the identification $Z_P\cong(\alpha\times\beta)\times Y_P$ of the pair
    effect $g$ and the identification of the configurations of
    $Y^{pq}_{\alpha^m,\beta^m}$ with $(\alpha\times\beta)^m\times Y_P$, the insertion
    word of $g$ at position $k$ has the operator $I_k\otimes1_{Y_P}$, where
    $I_k:\C^{\alpha\times\beta}\to(\C^{\alpha\times\beta})^{\otimes m}$ places its
    input in factor $k$ and $\eta$ in the other factors
    \cite[Lemma~5.1]{OpenAI2026PolynomialPEPS}.
\end{theorem}

\begin{proof}\leanok
    \uses{thm:party_op_prep_word, def:party_relabel, def:party_insertion_matrices}
    The operator is the permutation matrix of the relabelling times the operator
    $x\mapsto\eta^{\otimes(m-1)}\otimes x$ of Theorem~\ref{thm:party_op_prep_word}.
    The insertion $I_k$ is the preparation of $\eta^{\otimes(m-1)}$ behind its input
    followed by the relabelling $(i,f)\mapsto f'$, where $f'$ has $i$ at position $k$
    and $f$ at the other positions in order. Comparing entries at $p$, at $q$ and at
    the other parties proves the identity.
\end{proof}

\subsection{The stacks of a list of effects}
\label{sec:party_effect_stack_lists}

For a pair-source occurrence $e$ with endpoints $p_e\ne q_e$, fresh factors
$\alpha_e,\beta_e$ and vector $\eta_e$, write $T^m_e$ for $T$ with a stack of $m$ pair
registers on $p_e,q_e$ with pair basis $\alpha_e\times\beta_e$.

\begin{definition}[Output stacks of a list of effects]
    \label{def:party_effect_stack}
    \lean{QuantumCircuit.Party.stackApp, QuantumCircuit.Party.stackIdx,
        QuantumCircuit.Party.stackAppEquiv, QuantumCircuit.Party.stackIdxVector,
        QuantumCircuit.Party.norm_stackIdxVector, QuantumCircuit.Party.prepStack,
        QuantumCircuit.Party.frameStack, QuantumCircuit.Party.isAllowed_prepStack,
        QuantumCircuit.Party.sources_prepStack,
        QuantumCircuit.Party.isAllowed_frameStack,
        QuantumCircuit.Party.sources_frameStack}
    \leanok
    \uses{def:party_stack_system, def:party_prep_word, def:party_lift_pair,
        lem:kronecker_vectors}
    Let $E=(e_1,\dots,e_\ell)$ be a list of pair-source occurrences and $m\in\N$. The
    stacks of $E$ in front of $T$ form the system $T^m_E$, defined by $T^m_{()}=T$ and
    $T^m_{eE'}=(T^m_{E'})^m_e$, the first effect outermost. The stack
    configurations of $E$ are
    $J_m(E)=\prod_{i=1}^\ell(\alpha_{e_i}\times\beta_{e_i})^m$, in the order of the
    list, and the configurations of $T^m_E$ are identified with $T_P\times J_m(E)$.
    The ideal content of the stacks is
    \begin{align}
        \Gamma_m(E)
         & =\eta_{e_1}^{\otimes m}\otimes\cdots\otimes\eta_{e_\ell}^{\otimes m}
        \in\C^{J_m(E)},
        \label{eq:party_stack_content}
    \end{align}
    the unit vector of the one-point set for the empty list; it is a unit vector
    when every $\eta_{e_i}$ is. The preparation of the stacks is the word from $T$ to
    $T^m_E$ that prepares, for $E=eE'$, the stacks of $E'$ and then $m$ copies of
    $\eta_e$ on $p_e,q_e$. It is allowed when every $\eta_{e_i}$ is normalized, and
    its pair sources are $m$ copies of $(p_{e_1},q_{e_1},\eta_{e_1})$, then $m$
    copies of the next occurrence, and so on. A word $w$ from $S$ to $T$ beside the
    stacks of $E$ is the word from $S^m_E$ to $T^m_E$ obtained by placing $w$ beside
    the stack of every occurrence in turn; it is allowed when $w$ is and has the pair
    sources of $w$ \cite[Lemma~5.1]{OpenAI2026PolynomialPEPS}.
\end{definition}

\begin{theorem}[Operators of the stacks]
    \label{thm:party_op_effect_stack}
    \lean{QuantumCircuit.Party.op_prepStack, QuantumCircuit.Party.op_frameStack}
    \leanok
    \uses{def:party_effect_stack, def:party_effect_matrices}
    Read along the identification of the configurations of $T^m_E$ with
    $T_P\times J_m(E)$, the preparation of the stacks of $E$ has the operator
    $x\mapsto x\otimes\Gamma_m(E)$, and a word $w$ beside the stacks of $E$ has the
    operator $\mathrm{op}(w)\otimes1_{J_m(E)}$.
\end{theorem}

\begin{proof}\leanok
    \uses{thm:party_op_prep_word, thm:party_op_lift_pair}
    By induction on the list $E$. For $E=eE'$ the operator of the preparation is
    the operator $y\mapsto\eta_e^{\otimes m}\otimes y$ of
    Theorem~\ref{thm:party_op_prep_word} times $x\mapsto x\otimes\Gamma_m(E')$,
    and $\Gamma_m(eE')=\eta_e^{\otimes m}\otimes\Gamma_m(E')$. The operator of a
    word beside the stack of $e$ is $1\otimes\mathrm{op}(w')$ by
    Theorem~\ref{thm:party_op_lift_pair}, with $w'$ the word beside the stacks of
    $E'$; comparing entries gives both identities.
\end{proof}

\begin{theorem}[Ownership and dimensions of the stacks]
    \label{thm:party_stack_ownership}
    \lean{QuantumCircuit.Party.stackApp_of_forall_ne,
        QuantumCircuit.Party.card_stackApp}
    \leanok
    \uses{def:party_effect_stack}
    At a party $s$ that is no endpoint of an occurrence of $E$, the stacked system
    $T^m_E$ has the basis $T_s$. At every party $s$,
    \begin{align}
        \lvert(T^m_E)_s\rvert
         & =\lvert T_s\rvert\prod_{e\in E}
        \lvert\alpha_e\rvert^{m[s=p_e]}\,\lvert\beta_e\rvert^{m[s=q_e]},
        \notag
    \end{align}
    where $[s=p]$ is $1$ if $s=p$ and $0$ otherwise
    \cite[Lemma~5.1]{OpenAI2026PolynomialPEPS}.
\end{theorem}

\begin{proof}\leanok
    \uses{def:party_stack_system, def:party_add_pair}
    By induction on $E$: a stack changes no basis outside its endpoints, and it
    multiplies the dimension at $p_e$ by $\lvert\alpha_e\rvert^m$ and that at $q_e$
    by $\lvert\beta_e\rvert^m$.
\end{proof}

\subsection{The chain form of a circuit of pair sources and pair effects}
\label{sec:party_chain_form}

A circuit of pair sources and pair effects is a circuit whose gates are pair
sources and pair effects (Definition~\ref{def:party_pair_gates}); a monomial is
such a circuit on the parties of a cell. An effect chain
$F_nE_n\cdots E_1F_0$ is as in Definition~\ref{def:party_effect_chain}.

\begin{definition}[The chain form]
    \label{def:party_chain_form}
    \lean{QuantumCircuit.Party.ChainForm, QuantumCircuit.Party.ChainForm.ofWord,
        QuantumCircuit.Party.ChainCont, QuantumCircuit.Party.chainFormAux,
        QuantumCircuit.Party.chainFormAux_comp, QuantumCircuit.Party.chainFormAux_done,
        QuantumCircuit.Party.chainFormAux_localOp,
        QuantumCircuit.Party.chainFormAux_localOps,
        QuantumCircuit.Party.chainFormAux_source,
        QuantumCircuit.Party.chainFormAux_effect,
        QuantumCircuit.Party.PairEffectGate.toChain,
        QuantumCircuit.Party.Word.thenEffect,
        QuantumCircuit.Party.PairEffectGate.op_toChain,
        QuantumCircuit.Party.Monomial.chainForm,
        QuantumCircuit.Party.Monomial.toEffectChain,
        QuantumCircuit.Party.Monomial.chainEffects,
        QuantumCircuit.Party.Monomial.stackEquiv,
        QuantumCircuit.Party.Monomial.replaceWord}
    \leanok
    \uses{def:party_effect_chain, def:party_effect_stack, def:party_insertion_word,
        def:party_lift_pair, def:party_pair_gates, def:party_word}
    A chain form from $X$ to $W$ consists of an effect chain $C$ from $X_P$ to $W_P$,
    a list $E$ of pair-source occurrences (its effects), for every $m$ a bijection
    between the stack registers of $C$ and the stack configurations $J_m(E)$, and for
    every $m$ and every $\kappa:\Fin{n}\to\Fin{m}$, $n$ the number of effects of $C$,
    a word $R_{m,\kappa}$ from $X$ to $W^m_E$. The chain form of a word $w$ has the
    chain $\mathrm{op}(w)$ without effects and the word $w$ itself. A completion rule
    $k$ from $Y$ to $W$ assigns a chain form $k(v)$ from $X$ to $W$ to every word $v$
    from any system $X$ to $Y$.

    For a circuit $c$ of pair sources and pair effects from $Z$ to $Y$, a completion
    rule $k$ from $Y$ to $W$ and a word $w$ from $X$ to $Z$, the chain form
    $\Phi(c,k)(w)$ is defined by recursion on $c$
    \cite[Lemma~5.1]{OpenAI2026PolynomialPEPS}.
    \begin{itemize}
        \item The empty circuit gives $k(w)$; a local operation, local operations at
              every party, or a pair source $h$ give $k(h\circ w)$.
        \item A pair effect $g$ on $p\ne q$ with pair basis $\alpha\times\beta$ and
              vector $\eta$ closes the current piece: let $k(\emptyset)$ be the
              completion of the empty word on $Y$, with chain $C'$, effects $E'$ and
              words $R'$. The chain of $\Phi(c,k)(w)$ is $C'E_gF$, where $F$ is
              $\mathrm{op}(w)$ read along $Z_P\cong(\alpha\times\beta)\times Y_P$ and
              $E_g$ is the effect of $\eta$; its effects are $(p,q,\eta)$ followed by
              $E'$, its stack registers are identified factor by factor, and its
              words are $R_{m,\kappa}=\tilde R'_{m,\kappa'}\circ\iota_{\kappa_0}\circ
                  w$, where $\iota_{\kappa_0}$ is the insertion word of $g$ at
              position $\kappa_0$, $\kappa'$ is $\kappa$ without its first value, and
              $\tilde R'_{m,\kappa'}$ is $R'_{m,\kappa'}$ beside the new stack.
        \item A composition $c'\circ c$ gives $\Phi(c,\Phi(c',k))$.
    \end{itemize}
    One pair effect is a circuit of pair sources and pair effects with the operator
    of the pair effect, and for a word $w$, a pair effect $g$ and a circuit $c$ the
    circuit $c\circ g\circ w$ is the word, then the effect, then $c$.

    The chain form of a monomial $M$ from $X$ to $Y$ is $\Phi(M,k_0)(\emptyset)$,
    where $k_0$ assigns to a word its chain form. Its chain $C_M$, its effects $E_M$
    and its words $R^M_{m,\kappa}$, from $X$ to $Y^m_{E_M}$, are the effect chain,
    the effects and the replacement words of $M$. Composing with the bijection of
    stack registers identifies the configurations of $T^m_{E_M}$ with $T_P$ times
    the stack registers of $C_M$.
\end{definition}

\begin{theorem}[Invariants of the chain form]
    \label{thm:party_chain_form_invariants}
    \lean{QuantumCircuit.Party.eval_chainFormAux,
        QuantumCircuit.Party.isAllowed_chainFormAux,
        QuantumCircuit.Party.isAllowed_replaceWord_chainFormAux,
        QuantumCircuit.Party.stackVector_chainFormAux,
        QuantumCircuit.Party.op_replaceWord_chainFormAux,
        QuantumCircuit.Party.chainFormAux_toChain,
        QuantumCircuit.Party.Monomial.effects_chainFormAux,
        QuantumCircuit.Party.Monomial.effectCount_chainFormAux,
        QuantumCircuit.Party.Monomial.sourceCount_replaceWord_chainFormAux}
    \leanok
    \uses{def:party_chain_form, def:party_allowed, def:party_word,
        def:party_effect_chain_replacement}
    Let $c$ be a circuit of pair sources and pair effects from $Z$ to $Y$ and $k$ a
    completion rule from $Y$ to $W$. Each of the following holds for every word $w$
    into $Z$ when the corresponding hypothesis on $k$ holds for every word $v$ into
    $Y$.
    \begin{enumerate}
        \item If the chain of $k(v)$ evaluates to $R\,\mathrm{op}(v)$ for a fixed
              matrix $R$, the chain of $\Phi(c,k)(w)$ evaluates to
              $R\,\mathrm{op}(c)\,\mathrm{op}(w)$.
        \item Let $c$ be allowed. If the chain of $k(v)$ is allowed for allowed
              $v$, the chain of $\Phi(c,k)(w)$ is allowed for allowed $w$. If every
              word $R_{m,\kappa}$ of $k(v)$ is allowed for allowed $v$, every word
              $R_{m,\kappa}$ of $\Phi(c,k)(w)$ is allowed for allowed $w$.
        \item If the ideal content of the stacks of the chain of $k(v)$, read along
              its bijection of stack registers, is $\Gamma_m$ of its effects, the
              same holds for $\Phi(c,k)(w)$.
        \item If every word $R_{m,\kappa}$ of $k(v)$ has the operator of the term
              of the replaced chain of $k(v)$ with insertion positions $\kappa$, read
              along the identification of the configurations of $W^m_E$ with $W_P$
              times the stack registers, the same holds for $\Phi(c,k)(w)$.
        \item If $c$ is a word, then $\Phi(c,k)(w)=k(w')$ for a word $w'$ with
              $\mathrm{op}(w')=\mathrm{op}(c)\,\mathrm{op}(w)$, with the pair sources
              of $c$ followed by those of $w$, and allowed when $c$ and $w$ are.
        \item Let $c$ be a monomial. If the effects of $k(v)$ are a fixed list
              $E$, the effects of $\Phi(c,k)(w)$ are those of $c$ followed by $E$; if
              the chain of $k(v)$ has $n$ effects, the chain of $\Phi(c,k)(w)$ has
              those of $c$ plus $n$; if every $R_{m,\kappa}$ of $k(v)$ has as many
              sources as $v$ plus $N$, every $R_{m,\kappa}$ of $\Phi(c,k)(w)$ has
              as many sources as $w$, plus those of $c$, plus $m-1$ for every effect
              of $c$, plus $N$.
    \end{enumerate}
\end{theorem}

\begin{proof}\leanok
    \uses{def:party_word, def:party_pair_gates, thm:party_op_contraction,
        def:party_insertion_word, def:party_lift_pair, thm:party_op_insertion_word,
        thm:party_op_lift_pair, def:party_effect_chain,
        def:party_effect_chain_replacement, thm:party_word_op}
    Each statement is proved by induction on $c$, a composition $c'\circ c''$ using
    the statement for $c''$ completed by $\Phi(c',k)$. A local operation or a pair
    source joins the current word, whose operator is unchanged when it is read as a
    circuit of pair sources and pair effects. For a pair effect $g$, the effect
    matrix $E_g=\bra{\eta}\otimes1$ is the operator of $g$, so the chain evaluates
    to $R\,E_g\,\mathrm{op}(w)$; the piece $F$ is a contraction by
    Theorem~\ref{thm:party_op_contraction}; the insertion word and the word beside
    the stack are allowed, have $m-1$ and $N$ further sources, and have the
    operators $I_{\kappa_0}\otimes1$ and $1\otimes\mathrm{op}(R'_{m,\kappa'})$ by
    Theorems~\ref{thm:party_op_insertion_word} and~\ref{thm:party_op_lift_pair},
    which is the recursive formula of the replaced term after regrouping the stack
    configurations.
\end{proof}

\begin{theorem}[Monomials as effect chains]
    \label{thm:party_monomial_chain}
    \lean{QuantumCircuit.Party.Monomial.eval_toEffectChain,
        QuantumCircuit.Party.Monomial.isAllowed_toEffectChain,
        QuantumCircuit.Party.Monomial.chainEffects_eq,
        QuantumCircuit.Party.Monomial.effectCount_toEffectChain,
        QuantumCircuit.Party.Monomial.effectParties_eq,
        QuantumCircuit.Party.Monomial.norm_vector_of_mem_chainEffects,
        QuantumCircuit.Party.Monomial.toEffectChain_toChain,
        QuantumCircuit.Party.Monomial.toEffectChain_thenEffect,
        QuantumCircuit.Party.Monomial.chainEffects_thenEffect,
        QuantumCircuit.Party.Monomial.isAllowed_thenEffect}
    \leanok
    \uses{def:party_chain_form, def:party_allowed, def:party_pair_gates}
    For a monomial $M$ from $X$ to $Y$, the chain $C_M$ evaluates to $\mathrm{op}(M)$
    and has as many effects as $M$; the effects $E_M$ of the chain form are the pair
    effects of $M$ in chronological order, so their endpoints are those of the pair
    effects of $M$. If $M$ is allowed, then $C_M$ is allowed and every vector of
    $E_M$ is normalized. A word $w$ read as a monomial has the chain
    $\mathrm{op}(w)$. For a word $w$ from $X$ to $Z$, a pair effect $g$ from $Z$ to
    $Y$ with vector $\eta$ and a monomial $c$ from $Y$ to $W$, the chain of
    $c\circ g\circ w$ is $C_cE_gF$ with $F$ the operator of $w$ read along the
    configurations of $g$, its effects are $g$ followed by $E_c$, and the circuit
    $c\circ g\circ w$ is allowed when $w$ and $c$ are allowed and
    $\lVert\eta\rVert=1$
    \cite[Lemma~5.1]{OpenAI2026PolynomialPEPS}.
\end{theorem}

\begin{proof}\leanok
    \uses{thm:party_chain_form_invariants, thm:party_op_contraction,
        def:party_pair_gates}
    Apply Theorem~\ref{thm:party_chain_form_invariants} with the completion $k_0$,
    for which the chain of $k_0(v)$ is $\mathrm{op}(v)$, has no effect, and is
    allowed for allowed $v$ by Theorem~\ref{thm:party_op_contraction}: the first
    statement with $R=1$, the second, and the sixth with $E$ empty and $n=0$.
    Normalization of the effects of an allowed monomial is part of allowedness.
    The chains of a word and of $c\circ g\circ w$ follow from the fifth statement;
    the operator of a pair effect has norm at most $\lVert\eta\rVert$.
\end{proof}

\begin{theorem}[The replaced terms of a monomial are source-only]
    \label{thm:party_monomial_replaced_terms}
    \lean{QuantumCircuit.Party.Monomial.isAllowed_replaceWord,
        QuantumCircuit.Party.Monomial.op_replaceWord,
        QuantumCircuit.Party.Monomial.sourceCount_replaceWord,
        QuantumCircuit.Party.Monomial.isAllowed_prepStack,
        QuantumCircuit.Party.Monomial.op_prepStack,
        QuantumCircuit.Party.Monomial.stackApp_of_notMem_effectParties}
    \leanok
    \uses{def:party_chain_form, def:party_effect_chain_replacement,
        def:party_effect_stack}
    Let $M$ be a monomial from $X$ to $Y$ with $n$ pair effects and $s_M$ pair
    sources, let $m\in\N$ and $\kappa:\Fin{n}\to\Fin{m}$. Read along the
    identification of the configurations of $Y^m_{E_M}$ with $Y_P$ times the stack
    registers of $C_M$, the word $R^M_{m,\kappa}$ has as operator the term of the
    replaced chain $C_M$ with insertion positions $\kappa$. It has
    $s_M+n(m-1)$ pair sources, and it is allowed when $M$ is. The preparation of the
    stacks of $E_M$ in front of $T$ has the operator $x\mapsto x\otimes\Gamma$,
    where $\Gamma$ is the ideal content of the stacks of $C_M$; it is allowed when
    $M$ is. At a party that is no endpoint of a pair effect of $M$, the system
    $T^m_{E_M}$ has the basis of $T$ \cite[Lemma~5.1]{OpenAI2026PolynomialPEPS}.
\end{theorem}

\begin{proof}\leanok
    \uses{thm:party_chain_form_invariants, thm:party_monomial_chain,
        thm:party_op_effect_stack, thm:party_stack_ownership, def:party_effect_stack}
    For the completion $k_0$, the word of a chain without effects is $v$ itself, and
    the replaced term of the chain $\mathrm{op}(v)$ is $\mathrm{op}(v)$ followed by
    the preparation of the unit vector of the one-point set. The fourth, second and
    sixth statements of Theorem~\ref{thm:party_chain_form_invariants} give the
    operator, allowedness and the source count. The operator of the preparation
    follows from Theorem~\ref{thm:party_op_effect_stack} and the third statement,
    its allowedness from the normalization of $E_M$
    (Theorem~\ref{thm:party_monomial_chain}), and the ownership from
    Theorem~\ref{thm:party_stack_ownership}, since the endpoints of $E_M$ are those
    of the pair effects of $M$.
\end{proof}

\subsection{Elimination of pair effects for a weighted gate}
\label{sec:party_weighted_elimination}

Let $D$ be a cell and $L=((c_1,M_1),\dots,(c_K,M_K))$ a list of coefficients
$c_\xi\in\C$ and monomials $M_\xi$ from $A$ to $B$ on the parties of $D$, such as the
terms of a weighted gate (Definition~\ref{def:party_weighted_gate}). Gate
expansions, their inventory, their ideal vector, the replaced gate and its term
list are those of Definition~\ref{def:party_gate_expansion}.

\begin{definition}[The gate expansion of a weighted gate]
    \label{def:party_weighted_expansion}
    \lean{QuantumCircuit.Party.WeightedGate.toGate,
        QuantumCircuit.Party.WeightedGate.length_toGate,
        QuantumCircuit.Party.WeightedGate.map_norm_toGate,
        QuantumCircuit.Party.WeightedGate.mem_toGate,
        QuantumCircuit.Party.WeightedGate.gate_toGate,
        QuantumCircuit.Party.WeightedGate.ev_eq_gate_toGate,
        QuantumCircuit.Party.WeightedGate.isAllowed_toGate}
    \leanok
    \uses{def:party_gate_expansion, def:party_chain_form, def:party_weighted_gate}
    The gate expansion $\mathcal G_L$ of $L$ is the list of pairs $(c_\xi,C_{M_\xi})$.
    It has $K$ terms, with the absolute values $\lvert c_\xi\rvert$, every term is
    $(c,C_M)$ for a term $(c,M)$ of $L$, and its gate is
    $G=\sum_\xi c_\xi\,\mathrm{op}(M_\xi)$; in particular the matrix of a weighted
    gate is the gate of the expansion of its terms. When every $M_\xi$ is allowed,
    every chain of $\mathcal G_L$ is allowed.
\end{definition}

\begin{definition}[The output system of the replaced gate and its words]
    \label{def:party_replaced_gate_output}
    \lean{QuantumCircuit.Party.WeightedGate.gateOut,
        QuantumCircuit.Party.WeightedGate.gateEquiv,
        QuantumCircuit.Party.WeightedGate.prepGate,
        QuantumCircuit.Party.WeightedGate.isAllowed_prepGate,
        QuantumCircuit.Party.WeightedGate.wordList,
        QuantumCircuit.Party.WeightedGate.isAllowed_wordList}
    \leanok
    \uses{def:party_weighted_expansion, def:party_chain_form, def:party_effect_stack,
        def:party_gate_expansion}
    Let $m\in\N$. The output system $B^m_L$ of the replaced gate is $B$ for the empty
    list and $(B^m_{L'})^m_{E_M}$ for $L=(c,M)L'$: the stacks of all effect
    occurrences of all monomials in front of $B$. Its configurations are identified
    with $B_P$ times the inventory of $\mathcal G_L$, one monomial at a time by the
    identifications of Definition~\ref{def:party_chain_form}. The preparation word
    $\pi_L$ from $B$ to $B^m_L$ is empty for the empty list and, for $L=(c,M)L'$, is
    $\pi_{L'}$ followed by the preparation of the stacks of $E_M$; it is allowed when
    every monomial is. The word list of the empty list is empty. The word list of
    $L=(c,M)L'$ consists of the pairs of the first kind
    $(c\,m^{-n_M},\tilde\pi_{L'}\circ R^M_{m,\kappa})$, one for every
    $\kappa:\Fin{n_M}\to\Fin{m}$, with $n_M$ the number of effects of $C_M$ and
    $\tilde\pi_{L'}$ the word $\pi_{L'}$ beside the stacks of $E_M$, followed by the
    pairs of the second kind $(c',\rho\circ w')$, one for every term $(c',w')$ of the
    word list of $L'$, with $\rho$ the preparation of the stacks of $E_M$ in front of
    $B^m_{L'}$. When every monomial is allowed, every word of the list is allowed
    \cite[Lemma~5.1]{OpenAI2026PolynomialPEPS}.
\end{definition}

\begin{theorem}[The replaced gate on parties]
    \label{thm:party_replaced_gate_words}
    \lean{QuantumCircuit.Party.WeightedGate.op_prepGate,
        QuantumCircuit.Party.WeightedGate.termList_toGate,
        QuantumCircuit.Party.WeightedGate.isSourceOnly_termList,
        QuantumCircuit.Party.WeightedGate.gateOut_of_forall_ne,
        QuantumCircuit.Party.WeightedGate.card_gateOut}
    \leanok
    \uses{def:party_replaced_gate_output, def:party_gate_expansion, def:party_word}
    Read along the identification of the configurations of $B^m_L$ with $B_P$ times
    the inventory, the preparation word $\pi_L$ has the operator
    $x\mapsto x\otimes\Gamma_m$, with $\Gamma_m$ the ideal vector of $\mathcal G_L$,
    and the term list of the replaced gate of $\mathcal G_L$ is, term by term, the
    list of pairs $(c,\mathrm{op}(w))$ for the terms $(c,w)$ of the word list. When
    every monomial is allowed, every term of the term list, read on $B^m_L$, is
    source-only. At a party $s\in D$ that is no endpoint of a pair effect of any
    monomial of $L$, the basis of $B^m_L$ is $B_s$, and at every party $s$
    \begin{align}
        \lvert(B^m_L)_s\rvert
         & =\lvert B_s\rvert\prod_{\xi=1}^K\prod_{e}
        \lvert\alpha_e\rvert^{m[s=p_e]}\,\lvert\beta_e\rvert^{m[s=q_e]},
        \notag
    \end{align}
    the inner product over the pair effects $e$ of $M_\xi$
    \cite[Lemma~5.1]{OpenAI2026PolynomialPEPS}.
\end{theorem}

\begin{proof}\leanok
    \uses{thm:party_monomial_replaced_terms, thm:party_op_effect_stack,
        thm:party_stack_ownership, thm:party_monomial_chain,
        def:party_replaced_gate_output}
    By induction on $L$. For $L=(c,M)L'$ the operator of $\pi_L$ is the product of
    the operator $x\mapsto x\otimes\Gamma_{m}(C_M)$ of the preparation of the stacks
    of $E_M$ (Theorem~\ref{thm:party_monomial_replaced_terms}) and that of
    $\pi_{L'}$. A pair of the first kind has the
    operator $(\mathrm{op}(\pi_{L'})\otimes1)\,\mathrm{op}(R^M_{m,\kappa})$ by
    Theorem~\ref{thm:party_op_effect_stack}, which is the replaced term of $C_M$
    followed by the preparation of the ideal vector of $\mathcal G_{L'}$, the term of
    the term list; a pair of the second kind is a term of $L'$ followed by the
    preparation of the ideal content of the stacks of $C_M$. Allowed words give
    source-only terms. Ownership and dimensions follow from
    Theorems~\ref{thm:party_monomial_replaced_terms}
    and~\ref{thm:party_stack_ownership}, since the effects $E_M$ are the pair
    effects of $M$.
\end{proof}

\begin{theorem}[Elimination of normalized pair effects on parties]
    \label{thm:party_pair_effect_elimination}
    \lean{QuantumCircuit.Party.WeightedGate.partyPairEffectElimination}
    \leanok
    \uses{def:party_replaced_gate_output, def:party_weighted_expansion,
        def:party_gate_expansion, def:party_word}
    Let $m\ge1$ and $r\in\N$, and let every monomial $M_\xi$ of $L$ be allowed and have
    at most $r$ pair effects. Let $G=\sum_\xi c_\xi\,\mathrm{op}(M_\xi)$ be the gate
    of $\mathcal G_L$, let $G'_m$
    be the replaced gate of $\mathcal G_L$ and $\Gamma_m$ its ideal vector. Then
    $\lVert\Gamma_m\rVert=1$, the allowed word $\pi_L$ prepares $\Gamma_m$, read
    along the identification of the configurations of $B^m_L$ with $B_P$ times the
    inventory, and
    \begin{align}
        \lVert G'_m-G\otimes\ket{\Gamma_m}\rVert
         & \le\frac{r}{\sqrt m}\sum_{\xi=1}^K\lvert c_\xi\rvert.
        \label{eq:party_elimination_error}
    \end{align}
    Read along the same identification, $G'_m=\sum_{(c,w)}c\,\mathrm{op}(w)$ over
    the word list, every word of the list is allowed, and every term of the term
    list is source-only. The word list has at most $Km^r$ terms and the absolute sum
    $\sum_{(c,w)}\lvert c\rvert$ of its coefficients equals
    $\sum_\xi\lvert c_\xi\rvert$. At every party of $D$ that is no endpoint of a
    pair effect of any monomial, the output system $B^m_L$ of $G'_m$ has the basis of
    $B$.
\end{theorem}

This covers \cite[Lemma~5.1]{OpenAI2026PolynomialPEPS} on parties: the
error~(\ref{eq:party_elimination_error}), the expansion of the replaced gate into
local contractions and normalized pair sources, the bound $Km^r$ on the number of
monomials, the coefficient sum, here with equality, and the ownership of the
additional registers by the endpoint parties of the eliminated effects, which are
parties of the cell. The hypothesis that $G$ is a contraction is not used. The last
clause of the manuscript's lemma, that all sources on one pair of parties may be
combined into one normalized pair source, is not part of this theorem; for two pair
vectors the norm of the combined vector is Lemma~\ref{thm:peps_pair_combine}.

\begin{proof}\leanok
    \uses{thm:party_replace_gate, def:party_weighted_expansion,
        thm:party_monomial_chain, thm:party_replaced_gate_words,
        def:party_replaced_gate_output}
    Every chain of $\mathcal G_L$ is allowed and has as many effects as its monomial
    (Theorem~\ref{thm:party_monomial_chain}), so Theorem~\ref{thm:party_replace_gate}
    applies to $\mathcal G_L$: it gives the norm of $\Gamma_m$, the
    error~(\ref{eq:party_elimination_error}), the expansion of $G'_m$ into the term
    list, the bound on its length and the coefficient sum, since $\mathcal G_L$ has
    $K$ terms with the coefficients of $L$. Theorem~\ref{thm:party_replaced_gate_words}
    rewrites the term list as the word list and gives the preparation of
    $\Gamma_m$, source-only terms and ownership, and
    Definition~\ref{def:party_replaced_gate_output} gives that the words are
    allowed.
\end{proof}

\begin{definition}[Rescaling a weighted gate]
    \label{def:party_rescale}
    \lean{QuantumCircuit.Party.WeightedGate.rescale,
        QuantumCircuit.Party.WeightedGate.monomials_rescale,
        QuantumCircuit.Party.WeightedGate.sum_norm_rescale_le,
        QuantumCircuit.Party.WeightedGate.ev_rescale,
        QuantumCircuit.Party.WeightedGate.gate_toGate_rescale,
        QuantumCircuit.Party.WeightedGate.isAllowed_rescale}
    \leanok
    \uses{def:party_weighted_gate, def:party_weighted_expansion,
        lem:rescaled_contraction}
    For $\delta\in\R$, the rescaling of a weighted gate multiplies every coefficient
    by $(1+\delta)^{-1}$, with $0^{-1}=0$, and keeps the ordered list of monomials.
    Its matrix, and the gate of the expansion of its terms, are $(1+\delta)^{-1}$
    times those of the gate; it is allowed when the gate is, and for $\delta\ge0$
    its absolute coefficient sum is at most that of the gate. By
    Lemma~\ref{lem:rescaled_contraction}, a gate of norm at most $1+\delta$ thus
    becomes a contraction, as in the proof of
    \cite[Theorem~5.2]{OpenAI2026PolynomialPEPS}.
\end{definition}
